\FSseite{1 \quad Konstanten $\cdot$ Einheiten $\cdot$ Mathe $\cdot$ Fakten}{Bauelemente etit\,105}

\begin{multicols}{2}

% =====================================================================
\begin{FSblock}{1.1\ \ Konstanten \textnormal{— in Klausurgenauigkeit}}
\setlength{\tabcolsep}{1pt}
\begin{tabular}{@{}l@{\hspace{2mm}}l@{}}
$k_B$ & $=1{,}38\cdot10^{-23}\ \mathrm{J/K}$\\
$e$   & $=1{,}6\cdot10^{-19}\ \mathrm{As}$\\
$m_e$ & $=\erg{10^{-30}\,\mathrm{kg}}$ \warn{(Klausurvorgabe!)}\\
      & $=9{,}1\cdot10^{-31}\,\mathrm{kg}$ \ (Ü01)\\
$h$   & $=6{,}6\cdot10^{-34}\ \mathrm{Js}$\\
$\hbar$ & $=h/2\pi=1{,}05\cdot10^{-34}\ \mathrm{Js}$\\
$c_0$ & $=3\cdot10^{8}\ \mathrm{m/s}$\\
$\eps_0$ & $=8{,}854\cdot10^{-12}\ \mathrm{As/Vm}$\\
$\eps_{ox}$ & $=3{,}9$ \ (SiO$_2$)\\
$n_i$ & $=1{,}5\cdot10^{10}\,\mathrm{cm^{-3}}=1{,}5\cdot10^{16}\,\mathrm{m^{-3}}$\\
$W_g$ & $=1{,}1$\,eV (Si) \quad $0{,}67$\,eV (Ge)\\
$\mu_n$ & $=0{,}135\,\mathrm{m^2/Vs}=1350\,\mathrm{cm^2/Vs}$\\
$\mu_n{+}\mu_p$ & $=0{,}183\ \mathrm{m^2/Vs}$\\
$U_F$ & $=0{,}7$\,V (Si-Diode)\\
$U_{CE,sat}$ & $\approx0{,}2$\,V\\
\end{tabular}

\FSline
\textbf{Temperaturspannung} \ $\UT=\dfrac{k_B T}{e}$ \ bei $T=300$\,K:
\vspace{0.3mm}

\begin{tabular}{@{}l@{\hspace{1.5mm}}l@{}}
$\erg{25\,\mathrm{mV}}$ & \warn{Klausurvorgabe} — gilt in der Klausur\\
$25{,}9\,\mathrm{mV}$   & exakt aus $k_BT/e$\\
$26\,\mathrm{mV}$       & manche Übungen / Musterlösungen\\
\end{tabular}
\vspace{0.4mm}

\warn{\textbf{Achtung:}} $W_F-W_i=0{,}45$\,eV und $U_D=575$\,mV
entstehen nur mit $25{,}9\!-\!26$\,mV. Mit $25$\,mV kommt
$0{,}433$\,eV bzw. $555$\,mV heraus.
$\Rightarrow$ \erg{immer den vorgegebenen Wert einsetzen.}
\end{FSblock}

% =====================================================================
\begin{FSblock}{1.2\ \ Einheiten \warn{— häufigste Fehlerquelle}}
\centering
\setlength{\tabcolsep}{1pt}
\begin{tabular}{@{}l@{\,}l@{\hspace{3mm}}l@{\,}l@{\hspace{3mm}}l@{\,}l@{}}
G & $10^{9}$ & k & $10^{3}$   & $\mu$ & $10^{-6}$\\
M & $10^{6}$ & c & $10^{-2}$  & n & $10^{-9}$\\
  &          & m & $10^{-3}$  & p & $10^{-12}$\\
\end{tabular}
\justifying

\FSline
\textbf{Die drei Umrechnungen, die hier zählen}
\begin{align*}
1\,\mathrm{cm^{-3}} &= \erg{10^{6}}\,\mathrm{m^{-3}}
   \notat{$(10^{-2}\,\mathrm{m})^{-3}$}\\
1\,\mathrm{cm^{2}/Vs} &= 10^{-4}\,\mathrm{m^{2}/Vs}\\
1\,\mathrm{m^{2}/Vs} &= \erg{10^{4}}\,\mathrm{cm^{2}/Vs}
\end{align*}

\textbf{Merkfall MOSFET} (so in der Klausur):\\
$0{,}173\,\mathrm{m^2/Vs}\ \overset{\times10^4}{\longrightarrow}\
 \erg{1730\,\mathrm{cm^2/Vs}}$
\vspace{0.4mm}

\textbf{Merkfall Dotierung:}\\
$5\cdot10^{17}\,\mathrm{cm^{-3}}\ \overset{\times10^6}{\longrightarrow}\
 \erg{5\cdot10^{23}\,\mathrm{m^{-3}}}$
\vspace{0.4mm}

$1\,\mathrm{eV}=1{,}6\cdot10^{-19}\,\mathrm{J}$;\quad
$\mathrm{J}\to\mathrm{eV}$: \erg{$\div\,1{,}6\cdot10^{-19}$}
\tr{Zehnerpotenzen mit \textsf{EXP}/\textsf{EE} eingeben, nie als
$\times10\,\hat{}\,x$ — spart Klammerfehler.}
\end{FSblock}

% =====================================================================
\begin{FSblock}{1.3\ \ Mathe: Potenzen, Ableitung, Reihen}
$a^m a^n=a^{m+n}$ \quad $\dfrac{a^m}{a^n}=a^{m-n}$ \quad $(a^m)^n=a^{mn}$\\
$\sqrt[n]{a}=a^{1/n}$ \qquad $a^{-n}=\dfrac{1}{a^n}$

\FSline
$\dfrac{\dd}{\dd x}x^n=n\,x^{n-1}$ \quad
$\displaystyle\int x^n\dd x=\frac{x^{n+1}}{n+1}$ \quad
$\dfrac{\dd}{\dd x}\ee^{ax}=a\,\ee^{ax}$

\FSline
\textbf{Reihenentwicklung} \ (Oberwellen, S3)
\[ \ee^{x}=1+\frac{x}{1!}+\frac{x^{2}}{2!}+\frac{x^{3}}{3!}+\dots \]
\textbf{Additionstheorem} \ (Oberwellen, S3)
\[ 2\sin^{2}\omega t=1-\cos 2\omega t \]

\FSline
\textbf{Log.\ Differentiation} (Varaktor, S4)\\
$y=a\,x^{n}\ \Rightarrow\ \ln y=\ln a+n\ln x
 \notat{$\dd(\ )$}\ \dfrac{\Delta y}{y}=n\,\dfrac{\Delta x}{x}$

\FSline
\textbf{Mitternachtsformel} (MOSFET-AP, S6)
\[ ax^{2}+bx+c=0\ \Rightarrow\ x_{1,2}=\frac{-b\pm\sqrt{b^{2}-4ac}}{2a} \]
\tr{Beide Lösungen ausrechnen und \emph{beide} prüfen — in der
MOSFET-Aufgabe ist genau eine physikalisch möglich.}
\end{FSblock}

% =====================================================================
\begin{FSblock}{1.4\ \ Mathe: Netzwerk}
Reihe: $R=\sum R_i$,\quad $\dfrac{1}{C}=\sum\dfrac{1}{C_i}$\\
Parallel: $\dfrac{1}{R}=\sum\dfrac{1}{R_i}$,\quad $C=\sum C_i$\\
Zwei parallel: $R_1\!\parallel\!R_2=\dfrac{R_1R_2}{R_1+R_2}$\\
Spannungsteiler: $U_2=U\cdot\dfrac{R_2}{R_1+R_2}$\\
Maschenregel $\sum U=0$ \quad Knotenregel $\sum I=0$\\
Effektivwert: $\hat U=\sqrt{2}\,U_{e\!f\!f}$,\quad
$P=\dfrac{U_{e\!f\!f}^{2}}{R}$

\FSline
\textbf{Thévenin-Ersatzquelle} (grafische AP-Ermittlung, S3)
\begin{enumerate}
\item $U_0=$ Leerlaufspannung (Last abklemmen)
\item $R_i=$ Widerstand von den Klemmen aus,
      \warn{Spannungsquellen kurzschließen}
\end{enumerate}
\end{FSblock}

% =====================================================================
\begin{FSblock}{1.5\ \ Faktenblock \hfill \Quelle{Ü1 A3 · BA A1--A5 · PK-I A1 · PK13 A1 (jede Klausur)}}
\textbf{Licht — Spektrum}\\
Röntgen $0{,}01$–$10$\,nm $\cdot$ UV $10$–$400$\,nm $\cdot$
\erg{sichtbar $400$–$700$\,nm} $\cdot$ IR $700$\,nm–$1$\,mm $\cdot$
mm-Wellen $1$\,mm–$1$\,cm
\vspace{0.4mm}

\textbf{Zuordnung, die gefragt wird} (Wellenlänge $\to$ Bereich)\\
\setlength{\tabcolsep}{1.5pt}
\begin{tabular}{@{}l@{\ $\to$\ }l@{}}
$1$\,nm    & \erg{Röntgen} \ (PK-I A1)\\
$550$\,nm  & \erg{grün}, sichtbar \ (PK13 A1)\\
$700$\,nm  & rot, Grenze sichtbar/IR\\
$1{,}13\,\mu$m & IR — $\lambda_{max}$ von Si, s.\ S2\\
\end{tabular}

\FSline
\textbf{Größenordnungen}\\
Atomdurchmesser $\approx0{,}1$\,nm $\cdot$ $10^{23}$\,Atome/cm$^{3}$
$\cdot$ $10^{15}$\,Atome/cm$^{2}$ Oberfläche $\cdot$
„rein“ $=0{,}1\,\%$ Verunreinigung $\cdot$ Bindungsenergien $1$–$2$\,eV
$\cdot$ Aktivierungsenergie Dotanden $30$–$50$\,meV
(\erg{Klausurwert $35$\,meV}) $\cdot$ $k_BT|_{300\,\mathrm{K}}\approx25$\,meV
$\cdot$ $n_i$(Si) $=1{,}5\cdot10^{10}\,\mathrm{cm^{-3}}$
$\cdot$ $W_g$(Si) $=1{,}1$\,eV $\cdot$ $U_F$(Si) $=0{,}7$\,V

\FSline
\textbf{Dotierung}\\
n: \erg{P, As} — fünfwertig, geben Elektron ab\\
p: \erg{B, Al, In} — dreiwertig, nehmen Elektron auf\\
Si $=$ Grundmaterial (vierwertig)

\FSline
\textbf{Halbleiter} (Ankreuzfragen — Wortlaut merken)\\
$\sigma$ liegt \emph{zwischen} Metall und Isolator $\cdot$ $W_g$ kleiner
als beim Isolator $\cdot$
\erg{$\sigma$ nimmt mit \emph{fallender} Temperatur \emph{ab}} $\cdot$
$\sigma=$ Trägerdichte $\times$ Beweglichkeit $\times$ Ladung

\FSline
\textbf{Bauelemente / Begriffe in einem Satz}\\
Schottky-Diode $=$ Metall--Halbleiter-Übergang\\
CMOS $=$ Complementary Metal Oxide Semiconductor\\
Solarzelle $=$ innerer Fotoeffekt am pn-Übergang\\
Gleichrichten $=$ Strom nur in \emph{eine} Richtung durchlassen;
Wechsel- $\to$ Gleichgröße (Diode)\\
Halbleiter $=$ Festkörper mit $\sigma$ \emph{zwischen} Metall und Isolator,
$W_g$ klein genug, dass thermisch Paare entstehen\\
FET steuert \erg{leistungslos} über das E-Feld, BJT braucht
\erg{Steuerleistung} (Basisstrom)\\
Energieeigenwerte im unendlich hohen Topf: \erg{gequantelt}
(nicht kontinuierlich)\\
$W_g$ Isolator $\gg$ Halbleiter; $\sigma$ Metall $\gg$ HL $\gg$ Isolator
\end{FSblock}

% =====================================================================
\begin{FSblock}{1.6\ \ Definitionen in Antwortform \hfill \Quelle{Ü1 A2 · Ü4 A3 · PK13 A2}}
\setlength{\tabcolsep}{1.5pt}
\begin{tabularx}{\linewidth}{@{}>{\bfseries\RaggedRight}p{17mm}@{\hspace{1.3mm}}>{\RaggedRight\arraybackslash}X@{}}
Eigen-\newline leitung & Leitung im \emph{un}dotierten HL, allein durch thermisch
erzeugte Paare; $n=p=n_i$.\\[0.3mm]
therm.\ GGW & Paarbildung und Rekombination halten sich die Waage
$\Rightarrow$ $n\cdot p=n_i^{2}$.\\[0.3mm]
Paar-\newline bildung & Energie hebt ein Elektron von $W_V$ nach $W_C$;
es entsteht \emph{gleichzeitig} ein Loch.\\[0.3mm]
Rekom-\newline bination & Elektron fällt zurück ins Loch; Energie wird als Wärme
oder Photon frei.\\[0.3mm]
Dotieren & Gezieltes Einbringen von Fremdatomen, um \emph{eine}
Trägersorte um Größenordnungen zu erhöhen.\\[0.3mm]
n-/p-Leitung & $N_D>N_A$: Elektronen sind Majoritäten (n).
$N_A>N_D$: Löcher (p).\\[0.3mm]
Majori-\newline täten & häufigere Sorte ($n_0$ im n-HL); Minoritäten $=$ seltenere,
$p_0=n_i^{2}/n_0$.\\[0.3mm]
Störstellen-\newline erschöpfung & Bereich, in dem \emph{alle} Dotanden ionisiert
sind, Eigenleitung aber noch vernachlässigbar $\Rightarrow$
$n_0\approx N_D$ konstant. \erg{Arbeitsbereich.}\\[0.3mm]
Warum neutral? & Das abgegebene Elektron ist frei beweglich, der
zurückbleibende \emph{ortsfeste} Donatorrumpf positiv — in Summe null.\\[0.3mm]
Kontinui-\newline tätsgl. & Änderung der Trägerdichte $=$ Zufluss $-$ Abfluss
$+$ Generation $-$ Rekombination.\\[0.3mm]
Drift & Bewegung \emph{durch das E-Feld}: $j=\sigma E$.\\[0.3mm]
Diffusion & Bewegung \emph{durch Konzentrationsgefälle}:
$j=-e\,D\,\dd n/\dd x$.\\
\end{tabularx}
\end{FSblock}

\end{multicols}
